\section{Results}
\label{sec:results}

\subsection{Primary Result: erank--Oracle Correlation for BERT-base-uncased}

\begin{table}[h]
\centering
\caption{erank($W_0$) vs.\ PARA oracle rank correlation (BERT-base-uncased).}
\label{tab:primary}
\begin{tabular}{llll}
\toprule
Metric & Value & Threshold & Status \\
\midrule
Pearson $r$ (erank vs oracle rank) & \textbf{0.984} & $\geq 0.65$ & PASS \\
One-tailed $p$-value & \textbf{0.0013} & $< 0.05$ & PASS \\
PR--erank Pearson $\rho$ & \textbf{0.968} & $\geq 0.80$ & PASS \\
Oracle layers measured & 5 & $\geq 8$ (target) & PARTIAL \\
Families satisfying threshold & 1/3 & $\geq 2/3$ & PENDING \\
\bottomrule
\end{tabular}
\end{table}

Figure~\ref{fig:scatter} (scatter plot) shows erank($W_0$) vs.\ PARA oracle rank for the 5 measured BERT-base-uncased layers. Two attention output layers (\texttt{encoder.layer.0.attention.output.dense.weight}, \texttt{encoder.layer.6.attention.output.dense.weight}; erank $\sim 557$--$590$) and one attention query layer (\texttt{encoder.layer.3.attention.self.query.weight}; erank $\sim 556$) receive oracle rank $r^* = 4$; two FFN intermediate layers (\texttt{encoder.layer.8.intermediate.dense.weight}, \texttt{encoder.layer.10.intermediate.dense.weight}; erank $\sim 704$--$710$) receive oracle rank $r^* = 64$. erank cleanly separates these structural tiers.

\textbf{Bimodal oracle structure:} The oracle assigns exclusively $r^* = 4$ or $r^* = 64$ across the 5 measured layers. erank perfectly discriminates between these tiers with a structural scalar, requiring no task data.

\subsection{Metric Agreement: erank vs.\ Participation Ratio}

PR($W_0$) agrees with erank at $\rho = 0.968$, satisfying P5 (threshold $\geq 0.80$). Both structural metrics capture the same underlying signal. The one-tailed $p$-value of $0.0013$, well below $\alpha = 0.05$, provides statistical confirmation that the Pearson $r$ excludes zero without relying on bootstrap resampling. Bootstrap CI computation returned NaN values at $n=5$ due to the degenerate bimodal structure of the data (all resamples yield near-identical splits), and is noted for future reporting once full oracle coverage is achieved.

\subsection{erank Maps Across Model Families}

\begin{table}[h]
\centering
\caption{erank statistics across model families.$^{\dagger}$}
\label{tab:erank_stats}
\begin{tabular}{lccccc}
\toprule
Model & $n$ layers & erank min & erank max & Range ratio & CV \\
\midrule
BERT-base-uncased & 72+1 & 543.2 & 726.6 & 1.34$\times$ & $\sim$0.04 \\
DeBERTa-v3-base & 72 & 536.2 & 723.0 & 1.35$\times$ & $\sim$0.04 \\
ViT-base-patch16-224 & 73 & 244.4 & 727.2 & \textbf{2.97$\times$} & $\sim$0.12 \\
\bottomrule
\end{tabular}
\end{table}

$^{\dagger}$ BERT erank range excludes pooler.dense (erank = 405.6), which is not a LoRA adaptation target and is an architectural outlier; the range reported reflects LoRA-targetable layers only.

ViT-base shows 3$\times$ larger within-model erank variation than NLP encoders, suggesting stronger discriminative power if oracle ranks are available.

\subsection{Layer-Type Structure}

Figure~\ref{fig:heatmap} (erank heatmap, BERT) reveals consistent structure: FFN intermediate/output layers (erank $\sim 700$--$726$) consistently exceed attention Q/K/V/O layers (erank $\sim 530$--$611$). This layer-type hierarchy mirrors the oracle rank bimodality --- the oracle independently discovers the same structure erank predicts from $W_0$ alone.
